\section*{Interpreter Support}

\begin{itemize}
  \item \lstinline{apply_in_underlying_python(f, xs)}: \textit{primitive}, calls the function \lstinline{f}
    with arguments \lstinline{xs}. For example:
    \begin{lstlisting}
    def times(x, y):
        return x * y
    apply_in_underlying_python(times, llist(2, 3))  # returns 6
    \end{lstlisting}
  \item \lstinline{tokenize(x)}: \textit{primitive}, returns the literal linked list of tokens that results from tokenizing
    the string \lstinline{x} as a Python program. Each token is a string that contains the characters of
    the token as they appear in the program. Comments are ignored.
  \item \lstinline{parse(x)}: \textit{primitive}, returns the parse tree that results from parsing
    the string \lstinline{x} as a Python program. The following two pages describe the shape of the parse tree.
    The tree is represented by the tagged linked lists on the right; the angle brackets denote recursive application
    of the transformation rules. Implementations are allowed to support more of Python than listed.
\end{itemize}

In addition, the Source Academy frontend predeclares the name \lstinline{__program__} in all Python languages to
refer to the string representation of the entrypoint file of the program in the editor that is being run using ``Run''.
The entrypoint file is the file containing the code that acts as the entrypoint of the program being run.
If \lstinline{__program__} is used in the REPL, it refers to the string representation of the editor content
at the time when ``Run'' was last pressed.

\newpage
\KOMAoptions{paper=landscape,pagesize}
\recalctypearea
  \addtolength{\oddsidemargin}{-6cm}
  \addtolength{\evensidemargin}{-6cm}
  \addtolength{\textwidth}{13cm}

\begin{alignat*}{9}
  && \textit{program}      && ::= &\quad && \textit{statement}\ \ldots
                                                            && \texttt{llist("sequence", llist of  } \langle \textit{statement}\rangle \texttt{)} \textrm{; a single statement stands for itself}                                                                              \\[1mm]
  && \textit{statement}    && ::= &\quad && \textit{name}\ \textbf{\texttt{=}}\  \textit{expression}
                                                            && \texttt{llist("assignment",  } \langle \textit{name}\rangle \textit{,  } \langle \textit{expression}\rangle \texttt{)}                                      \\
  &&                       && |   &\quad && \textbf{\texttt{def}}\  \textit{name} \ \textbf{\texttt{(}}\  \textit{names} \ \textbf{\texttt{)}}\ \textit{: body} \quad
                                                            &&  \texttt{llist("function\_definition",  } \langle \textit{name}\rangle \textit{, } \langle \textit{names}\rangle \textit{,}\ \langle \textit{body}\rangle \texttt{)} \\
  &&                       && |   &\quad && \textbf{\texttt{return}}\  \textit{expression}
                                                            && \texttt{llist("return\_statement",  } \langle \textit{expression}\rangle \texttt{)}                                                                                    \\
  &&                       && |   &\quad && \textit{if-statement} \quad
                                                            && \textrm{see below}                                                                                                                                                            \\
  &&                       && |   &\quad && \textbf{\texttt{while}}\ \textbf{\texttt{(}}\  \textit{expression} \ \textbf{\texttt{)}} \ \textit{: body}
                                                            && \texttt{llist("while\_loop",  } \langle \textit{expression}\rangle \textit{,  } \langle \textit{body}\rangle \texttt{)}                                               \\
  &&                       && |   &\quad && \textbf{\texttt{for}}\ \textit{name}\ \textbf{\texttt{in range}}\ \textit{( range-args ) : body}
                                                            && \texttt{llist("for\_loop",  } \langle \textit{name}\rangle \textit{,  } \langle \textit{range\_args}\rangle \textit{,  } \langle \textit{body}\rangle \texttt{)}                                                \\
  &&                       && |   &\quad && \textbf{\texttt{break}}
                                                            && \texttt{llist("break\_statement")}                                                                                                                                     \\
  &&                       && |   &\quad && \textbf{\texttt{continue}}
                                                            && \texttt{llist("continue\_statement")}                                                                                                                                  \\
  &&                       && |   &\quad && \textbf{\texttt{global}}\ \textit{names}
                                                            && \texttt{llist("global\_statement",  } \langle \textit{names}\rangle \texttt{)}                                                                                          \\
  &&                       && |   &\quad && \textbf{\texttt{nonlocal}}\ \textit{names}
                                                            && \texttt{llist("nonlocal\_statement",  } \langle \textit{names}\rangle \texttt{)}                                                                                        \\
  &&                       && |   &\quad && \textbf{\texttt{pass}}
                                                            && \texttt{llist("pass\_statement")}                                                                                                                                      \\
  &&                       && |   &\quad && \textit{expression}
                                                            && \textrm{see below}                                                                                                                                                            \\[1mm]
  && \textit{range\_args}  && ::= &\quad && \textit{expression}_1 \textbf{\texttt{[}} \textit{, expression}_2 \textbf{\texttt{]}} \textbf{\texttt{[}} \ \textit{, expression}_3 \textbf{\texttt{]}} \
                                                            &&  \texttt{llist("range\_args",  } \langle \textit{expression}_1\rangle \textit{, } \langle \textit{expression}_2\rangle \textit{,}\ \langle \textit{expression}_3\rangle \texttt{)} \\
  && \textit{names}        && ::= &\quad && \epsilon\ | \  \textit{name} \ (\ \textbf{\texttt{,}} \ \textit{name}\ )\ \ldots
                                                            && \texttt{llist of  } \langle \textit{name}\rangle                                                                                                                       \\
  && \textit{if-statement} && ::= &\quad && \term{if}\ \textit{expression}_1\ \term{:}\ \textit{body}_1 \\
  &&                       &&     &\quad && (\ \term{elif}\ \textit{expression}_2\ \term{:}\ \textit{body}_2\ )...                                                                                                                          \\
  &&                       &&     &\quad && \term{else}\ \term{:}\ \textit{body}_3
                                                            && \texttt{llist("conditional\_statement",  } \langle \textit{expression}_1\rangle \textit{, }                                                                            \\
                                                            &&&&&&&&&\ \ \ \ \ \ \texttt{ } \langle \textit{body}_1\rangle \textit{,  } \langle \textit{body}_3\rangle \ \textrm{or}\ \langle \textit{if-statement} \rangle \texttt{)} \textrm{; a missing else part is } \texttt{llist("pass\_statement")}     \\
  && \textit{body}         && ::= &\quad && \textit{statement}... \quad
                                                            && \texttt{llist("sequence", llist of  } \langle \textit{statement}\rangle \texttt{)} \textrm{; a single statement stands for itself}                                                                                                 \\
\end{alignat*}

\begin{alignat*}{9}
  && \textit{expression}.  && ::= &\quad && \textit{integer} 
                                                            && \texttt{llist("literal",}\ \textit{integer}\texttt{)}     \\
  &&                       && |   &\quad && \textit{float}
                                                            && \texttt{llist("literal",}\ \textit{float}\texttt{)}       \\
  &&                       && |   &\quad && \textit{complex}
                                                            && \texttt{llist("literal",}\ \textit{complex}\texttt{)}     \\
  &&                       && |   &\quad && \term{True}\ |\ \term{False}
                                                            && \texttt{llist("literal",}\ \textit{True}\texttt{)}        \\
  &&                       && |   &\quad && \term{None}
                                                            && \texttt{llist("literal",}\ \textit{None}\texttt{)}        \\
  &&                       && |   &\quad && \textit{string}
                                                            && \texttt{llist("literal",}\ \textit{string}\texttt{)}      \\
  &&                       && |   &\quad && \textit{name}
                                                            && \texttt{llist("name", string)}                            \\
  &&                       && |   &\quad && \textit{expression}_1 \  \textit{log-op} \ \textit{expression}_2 \qquad
                                                            && \texttt{llist("logical\_composition",  } \langle \textit{log-op}\rangle \textit{, } \langle \textit{expression}_1\rangle \textit{, } \langle \textit{expression}_2\rangle \texttt{)}          \\
  &&                       && |   &\quad && \textit{expression}_1 \  \textit{bin-op} \ \textit{expression}_2 \qquad
                                                            && \texttt{llist("binary\_operator\_combination",  } \langle \textit{bin-op}\rangle \textit{, } \langle \textit{expression}_1\rangle \textit{, } \langle \textit{expression}_2\rangle \texttt{)} \\
  &&                       && |   &\quad && \textit{un-op} \ \textit{expression}
                                                            && \texttt{llist("unary\_operator\_combination",  } \langle \textit{un-op}\rangle \textit{, } \langle \textit{expression}\rangle \texttt{)}                                                      \\
  &&                       && |   &\quad && \textit{expression} \ \textbf{\texttt{(}}\ \textit{expressions}\ \textbf{\texttt{)}}
                                                            && \texttt{llist("application",  } \langle \textit{expression}\rangle \textit{, llist of  } \langle \textit{expression}\rangle \texttt{)}                                                 \\
  &&                       && |   &\quad && \textbf{\texttt{lambda}}\ \textit{names}\ \texttt{\textbf{:}}\ \textit{expression}
                                                            && \texttt{llist("lambda\_expression",  } \langle \textit{names}\rangle \textit{,  } \langle \textit{expression}\rangle \texttt{)}                                                                     \\
  &&                       && |   &\quad && \textit{expression}_1 \ \textbf{\texttt{if}}\ \textit{expression}_2 \ \textbf{\texttt{else}}\ \textit{expression}_3\
                                                            && \texttt{llist("conditional\_expression",  } \langle \textit{expression}_1\rangle \textit{,}                                                                                                   \\
                                                            &&&&&&&&&\ \ \ \ \ \ \texttt{ } \langle \textit{expression}_2\rangle \textit{,  } \langle \textit{expression}_3\rangle \texttt{)}                                                                       \\
  &&                       && |   &\quad && \textit{expression}_1 \textbf{\texttt{[}}
                                            \textit{expression}_2 \textbf{\texttt{]=}}\ \textit{expression}_3
                                                            && \texttt{llist("list\_assignment",  } \langle \textit{expression}_1\rangle \textit{,  } \langle \textit{expression}_2\rangle \textit{,  } \langle \textit{expression}_3\rangle \texttt{)}      \\
  &&                       && |   &\quad && \textit{expression}_1 \textbf{\texttt{[}} \textit{expression}_2 \textbf{\texttt{]}}
                                                            && \texttt{llist("list\_access",  } \langle \textit{expression}_1\rangle \textit{,  } \langle \textit{expression}_2\rangle \texttt{)}                                                            \\
  &&                       && |   &\quad && \textbf{\texttt{[}}\ \textit{expressions}\ \textbf{\texttt{]}}
                                                            && \texttt{llist("list\_expression", llist of  } \langle \textit{expression}\rangle \texttt{)}                                                                                            \\
  &&                       && |   &\quad && \textbf{\texttt{(}}\  \textit{expression} \ \textbf{\texttt{)}} && \textrm{parenthesised}\ \textit{expression}                                            \\[1mm]
  && \textit{log-op}.      && ::= &\quad && \textbf{\texttt{and}}\ |\ \texttt{\textbf{or}}
                                                            && \textit{string representing operator}                                                                                                  \\[1mm]
  && \textit{bin-op}       && ::= &\quad && \textbf{\texttt{+}}\ |\ \textbf{\texttt{-}}\ |\ \textbf{\texttt{*}}\ |\ \textbf{\texttt{/}}\ |\ \textbf{\texttt{//}}\ |\ \textbf{\texttt{**}}\ |\ \textbf{\texttt{\%}}\ |\ \textbf{\texttt{is}}\ |\ \textbf{\texttt{is not}}
                                                            && \textit{string representing operator}                                                                                                  \\[1mm]
  &&                       && |   &\quad &&  \textbf{\texttt{== }} |\ \texttt{\textbf{!=}}\ |\ \texttt{\textbf{>}}\ |\ \texttt{\textbf{<}}\ |\ \texttt{\textbf{>= }} |\ \texttt{\textbf{<=}}
                                                            && \textit{string representing operator}                                                                                                  \\[1mm]
  && \textit{un-op}        && ::= &\quad && \textbf{\texttt{not }} |\ \textbf{\texttt{-}}
                                                            && \texttt{"-unary"}                                                                                                                      \\
  &&                       && |   &\quad && \textbf{\texttt{+}}
                                                            && \texttt{"+unary"}                                                                                                                      \\
  && \textit{expressions}  && ::= &\quad && \epsilon\ | \ \textit{expression}\ ( \ \textbf{\texttt{,}} \ \textit{expression} \ )\ \ldots
                                                            && \texttt{llist of  } \langle \textit{expression}\rangle                                                                          \\
\end{alignat*}



\newpage
\KOMAoptions{paper=portrait}
\recalctypearea

% ALIGN EVEN- AND ODD-NUMBERED PAGES.
\evensidemargin 35pt

% HORIZONTAL MARGINS
% Left margin 1 inch (0 + 1)
\setlength{\oddsidemargin}{0in}
% Text width 6.5 inch (so right margin 1 inch).
\setlength{\textwidth}{6.5in}

% ----------------
% VERTICAL MARGINS
% Top margin 0.5 inch (-0.5 + 1)
\setlength{\topmargin}{-0.5in}
% Head height 0.25 inch (where page headers go)
\setlength{\headheight}{0.25in}
% Head separation 0.25 inch (between header and top line of text)
\setlength{\headsep}{0.25in}
% Text height 8.5 inch (so bottom margin 1.5 in)
\setlength{\textheight}{10.0in}

